\chapter{Background and Related Work}\label{ch:background}

\begin{guide}
About 6--8 pages. This is the chapter you can reuse most: the proposal's ``Background and Related Work'' and your literature review already cover it. Paste it in from the proposal's \texttt{.tex} source, then:
\begin{itemize}
  \item cut anything about methods you ended up not using;
  \item add the short sections marked NEW below, because the results now depend on them;
  \item end the chapter with a ``research gap'' paragraph that leads straight into your hypotheses.
\end{itemize}
Copy the proposal's \texttt{.bib} entries into \texttt{references.bib}.
\end{guide}

\section{Deep learning for financial time series}
\tbc{from proposal: RNNs, vanishing gradients, LSTM, GRU, Bi-LSTM}

\section{Attention mechanisms in sequence modelling}
\tbc{from proposal}

\section{Multivariate features and cross-asset signals}
\tbc{from proposal: macro, on-chain, correlations}

\section{Volatility modelling and jump risk}
\tbc{from proposal: realised volatility, GARCH(1,1), jump-diffusion}

\section{Forecast evaluation and naive baselines}\label{sec:bg-baselines}
\begin{guide}
NEW. Explain why any forecast must be compared with a naive baseline: for returns, ``predict no change'' (random walk); for volatility, persistence (``tomorrow looks like today''). Cite the efficient-market / random-walk view of returns and the volatility-clustering literature. This sets up two of your main findings.
\end{guide}
\tbc{}

\section{Risk management and backtesting for leveraged trading}\label{sec:bg-risk}
\begin{guide}
NEW (or expand). Volatility targeting / volatility-scaled position sizing, ATR-style stops, leverage and margin, why drawdown and Sharpe are the standard measures, and the pitfalls of backtests (look-ahead bias, overfitting the tuning, costs). Mention EU retail leverage caps on crypto CFDs (2:1) and that the UK bans retail crypto CFDs - cite the regulators.
\end{guide}
\tbc{}

\section{Research gap}
\tbc{}
